\documentclass[final,3p,times]{elsarticle}

\usepackage{amsmath,amssymb,amsthm}
\usepackage[hidelinks]{hyperref}
\setlength{\bibsep}{0pt}
\biboptions{sort&compress}
\geometry{textheight=712pt,headheight=12pt}
\emergencystretch=1.5em

\newtheorem*{theorem}{Theorem}
\journal{Information Processing Letters}
\myfooter[L]{Intended for publication in Information Processing Letters\hfill\today}

\begin{document}

\begin{frontmatter}

\title{The Ford--Johnson algorithm, no longer unbeaten\\ for less than 47 elements}

\author{H. G. Sager}
\ead{hazim.g.sager@gmail.com}
\address{Independent researcher}

\begin{abstract}
It is shown that 26 elements can be sorted with 90 comparisons, one fewer than merge
insertion requires.
\end{abstract}

\begin{keyword}
Analysis of algorithms \sep Sorting \sep Merge insertion \sep Formal verification
\end{keyword}

\end{frontmatter}

In 2007 Peczarski~\cite{Pec07} found the algorithm of Ford and Johnson~\cite{FJ59}
`still unbeaten for less than 47 elements'. We show that it no longer is.

Let $S(n)$ denote the least number of comparisons sufficient to sort $n$ distinct
elements. The algorithm of Ford and Johnson, now usually called merge insertion, sorts $n$
elements with at most
\[
  F(n) = \sum_{k=1}^{n} \Bigl\lceil \log_2 \frac{3k}{4} \Bigr\rceil
\]
comparisons~\cite{Knuth98}, and it is known that $S(n) = F(n)$ for every $n \le
22$~\cite{Knuth98,Pec04,Pec07,SW23}. Manacher~\cite{Man79} proved that $S(n) < F(n)$ for
infinitely many $n$; the least such $n$ previously known was 47, by a result of Schulte
M\"onting~\cite{SM81}.

\begin{theorem}
$S(26) \le 90$ and $S(46) \le 194$, whereas $F(26) = 91$ and $F(46) = 195$. More
generally, $S(n) < F(n)$ whenever $26 \le n \le 32$ or $46 \le n \le 64$.
\end{theorem}

\begin{proof}
We sort 21 of the 26 elements by merge insertion, which requires $F(21) = 66$
comparisons. The remaining five have not yet been compared with anything, and they may be
inserted into the ordered chain of 21 by means of a decision tree $T_{21,5}$ of depth 24,
given in~\cite{data}. This makes 90 comparisons in all. In the same way, 42 of 46 elements
may be sorted with $F(42) = 171$ comparisons and the remaining four inserted by a tree
$T_{42,4}$ of depth 23, which makes 194. Finally, binary insertion gives $S(n+1) \le S(n) + \lceil \log_2 (n+1) \rceil$, and for
$26 \le n < 32$ and $46 \le n < 64$ this increment, five and six respectively, equals
$F(n+1) - F(n)$. Hence $S(n) \le F(n) - 1$ implies $S(n+1) \le F(n+1) - 1$, and the
improvement extends to 32 and 64 elements.
\end{proof}

Inserting the new elements one at a time by binary insertion would take 25 and 24
comparisons, and the information-theoretic lower bounds are 23 and 22; every leaf of either
tree is reached by exactly one input. Since $26! > 2^{88}$, at least 89 comparisons are
necessary, so $S(26)$ is either 89 or 90. For 47 elements the argument gives 200, the bound
of Schulte M\"onting. The proof has been formalised in Lean~4~\cite{Lean4}, using
\texttt{native\_decide} to check the trees because of their size; the formalisation and the
trees are available at~\cite{data}. They were found by computer search, the specific
technique of which may be described elsewhere.

\section*{Declaration of generative AI and AI-assisted technologies in the manuscript
preparation process}
During the preparation of this work the author used Claude Opus 5.5 (Anthropic) in order to
assist in writing up this note. After using this tool, the author reviewed and edited the
content as needed and takes full responsibility for the content of the publication.

\begin{thebibliography}{10}
\footnotesize
\setlength{\itemsep}{0pt}\setlength{\parskip}{0pt}

\bibitem{Pec07}
M. Peczarski, The Ford--Johnson algorithm still unbeaten for less than 47 elements,
Inform. Process. Lett. 101 (2007) 126--128.

\bibitem{FJ59}
L.R. Ford Jr., S.M. Johnson, A tournament problem, Amer. Math. Monthly 66 (1959) 387--389.

\bibitem{Knuth98}
D.E. Knuth, The Art of Computer Programming, Vol.~3: Sorting and Searching, 2nd ed.,
Addison-Wesley, 1998, Section 5.3.1.

\bibitem{Pec04}
M. Peczarski, New results in minimum-comparison sorting, Algorithmica 40 (2004) 133--145.

\bibitem{SW23}
F. Stober, A. Wei{\ss}, Lower bounds for sorting 16, 17, and 18 elements, in: Proc.
Symposium on Algorithm Engineering and Experiments (ALENEX), SIAM, 2023.
arXiv:2206.05597.

\bibitem{Man79}
G.K. Manacher, The Ford--Johnson sorting algorithm is not optimal, J. ACM 26 (1979)
441--456.

\bibitem{SM81}
J. Schulte M\"onting, Merging of 4 or 5 elements with $n$ elements, Theoret. Comput. Sci.
14 (1981) 19--37.

\bibitem{data}
H.G. Sager, The Ford--Johnson algorithm, no longer unbeaten for less than 47 elements,
v0.1.1, Zenodo, 2026, \url{https://doi.org/10.5281/zenodo.23663357}, source at
\url{https://github.com/softlytyped/ford-johnson-beaten}.

\bibitem{Lean4}
L. de Moura, S. Ullrich, The Lean 4 theorem prover and programming language, in: Automated
Deduction -- CADE 28, LNCS 12699, Springer, 2021, pp. 625--635.

\end{thebibliography}

\end{document}
